% varanc-presence.tex — Mixture-of-trees variational TKF92 ancestral presence/absence
% To be \input from tkf.tex

\section{Mixture-of-trees variational ancestral presence/absence}
\label{sec:varanc-presence}

This appendix derives a variational lower bound on the indel
log-likelihood of an MSA under the TKF92 conditional WFST
$\tkfwfst'(\insrate,\delrate,\evoltime,\ext)$ of
Section~\ref{sec:tkf92-wfst} when applied independently
on each branch of a phylogenetic tree, combined with a
TKF stationary HMM at the root.
The bound is parameterised by a tree-structured graphical model
on per-column ancestral presence/absence indicators with three
irreversible states \{NotYetInserted, Present, Deleted\}.
It is a proper lower bound on the full TKF92 marginal
log-likelihood $\log p(\text{MSA} \mid \parsetree, \theta)$, with
gap split into a standard $q$-dependent variational
KL term and a $q$-independent ``restriction'' term arising because
our latent space cannot represent ancestor residues that get
inserted and then fully deleted on every leaf-bound lineage before
reaching any leaf
(Section~\ref{sec:varanc-presence-restricted}).
For ancestral-presence inference at the observed columns the
restriction term is irrelevant, since it cancels from the q-optimal
posterior;
for parameter learning it does not, and a tighter bound would
require explicit modelling of those ghost-column histories.
Special-case parameter settings recover an irreversible 3-state
analogue of Felsenstein gap-CTMC reconstruction and Fitch parsimony.

\subsection{Setting and approximation}
\label{sec:varanc-presence-setting}

Fix a rooted phylogenetic tree $\parsetree$ with internal-node set
$\mathcal{I}$, leaf set $\mathcal{L}$, and a branch length
$\branchlen_e > 0$ on each edge $e$.
We are given a multiple-sequence alignment with $L$ columns whose rows
are the leaf sequences of $\parsetree$;
let $X^v_n \in \{0,1\}$ denote the observed presence indicator
($1$ if leaf $v \in \mathcal{L}$ has a residue in column $n$).
The task is to infer per-internal-node presence indicators
$\{X^v_n\}_{v \in \mathcal{I}, 1 \le n \le L}$.

We replace the full TKF92 generative model on the tree with the
following \emph{per-branch conditional approximation}:
each branch $v \to w$ is taken to evolve its row according to the
TKF92 conditional WFST $\tkfwfst'$ acting independently on the whole
parent sequence, with no shared latent fragment structure across
branches.
A residue inserted by the WFST on one branch may therefore be
partially deleted, character by character, on a later branch
(unlike in the strict TKF92 process, where fragments are indivisible).
The WFST $\tkfwfst'$ is itself a closed-form approximation to the
finite-time transition probabilities of the General Geometric Indel
(GGI) model — the maximum-entropy CTMC with given expected per-site
indel rates and expected inserted-fragment length — and has been
shown empirically to fit GGI Gillespie simulations
\cite{Holmes2020} and pairwise-alignment data
\cite{LargeHolmes2026} well.
Using it as the per-branch transition law of the tree-level model in
this appendix is therefore a small additional approximation on top of
the GGI approximation itself.

\subsection{Restricted generative model}
\label{sec:varanc-presence-restricted}

The model joint over leaf and internal-node presence patterns at the
$L$ MSA columns is
\begin{equation}
\label{eq:joint}
p(\text{MSA},\, \{X^v\}_{v \in \mathcal{I}} \mid \parsetree, \theta)
= p_{\text{singlet}}\!\big(X^{\text{root}} \mid \theta\big)\,
   \prod_{(v \to w) \in \parsetree}
   P\!\big(X^w \mid X^v, \branchlen_{vw}, \theta\big),
\end{equation}
where $p_{\text{singlet}}$ is the TKF stationary distribution
(Section~\ref{sec:varanc-presence-singlet}, below) and each
$P(X^w \mid X^v, t)$ is the per-branch presence-pattern conditional
defined by the WFST path log-likelihood
(Section~\ref{sec:varanc-presence-path-LL}).
Marginalising~\eqref{eq:joint} over internal patterns supported on
the $L$ observed MSA columns gives a quantity we denote
$\tilde{p}(\text{MSA} \mid \parsetree, \theta)$.
The full TKF92 marginal $p(\text{MSA} \mid \parsetree, \theta)$ also
sums over evolutionary histories with ancestral residues that get
inserted and then fully deleted on every leaf-bound lineage before
reaching any leaf — \emph{ghost columns} that leave no MSA trace
and so are not representable as values of $X^v_n$ at observed
columns $1, \ldots, L$.
Since each ghost-column history adds a non-negative probability to
$p$ that is absent from $\tilde{p}$,
\begin{equation}
\label{eq:tildep-le-p}
\tilde{p}(\text{MSA} \mid \parsetree, \theta)
\;\leq\;
p(\text{MSA} \mid \parsetree, \theta).
\end{equation}
The gap $\log p - \log \tilde{p} \geq 0$ depends on the model
parameters and the tree but not on the variational distribution $q$
introduced below.
For the ancestral-presence reconstruction task — fix $\theta$,
optimise $q$ to recover internal-node marginals at the observed
columns — this gap is a $q$-independent constant and the
restricted-model posterior at observed columns coincides with the
full-model posterior conditioned on the same support, so all
inferences below proceed unchanged.

\subsection{Singlet HMM at the root}
\label{sec:varanc-presence-singlet}

The TKF92 stationary HMM (Section~\ref{sec:tkf92-machines}) emits a
sequence of length $L_{\text{root}}$ with
\begin{equation}
\label{eq:singlet-tkf92}
p_{\text{singlet}}(L_{\text{root}})
= \begin{cases}
1 - \kappa & L_{\text{root}} = 0, \\
\kappa\, p^{L_{\text{root}} - 1}\,(1-\ext)(1-\kappa)
       & L_{\text{root}} \geq 1,
\end{cases}
\end{equation}
where $\kappa = \insrate/\delrate$ and $p = \ext + (1-\ext)\kappa$.
Identifying $L_{\text{root}}$ with the number of MSA columns at which
the root is in state $\state{P}$
($L_{\text{root}} = \sum_n \delta(X^{\text{root}}_n = 1)$) gives the
root prior in the joint~\eqref{eq:joint}.
The two-cases form arises because the TKF92 singlet HMM has distinct
``from $\sta$'' (probability $\kappa$ to emit, $1-\kappa$ to end) and
``from emitting state'' (probability $p$ to extend, $(1-\ext)(1-\kappa)$
to end) transition contexts.

For TKF91 ($\ext = 0$): $p = \kappa$ and the two cases collapse to the
single form $p_{\text{singlet}}(L_{\text{root}}) = (1-\kappa)\kappa^{L_{\text{root}}}$.

\paragraph{Note on a tempting shortcut.}
One might hope to avoid the singlet by interpreting the root as the
output of a pseudo-branch with $\evoltime \to \infty$ and an
all-absent pseudo-parent.
For TKF91 this works — the $t \to \infty$ limit of $\tkfwfst$ with an
empty input does recover the singlet sequence-length distribution.
For TKF92 it does \emph{not}: the WFST's $(1-\ext)$ factors are tied
to the fragment-extension structure of the Pair HMM and the limiting
output mass for length $L \geq 1$ is $\kappa\,p^{L-1}(1-\kappa)$
rather than the correct $\kappa\,p^{L-1}(1-\ext)(1-\kappa)$ —
i.e.\ the WFST does not normalise on an empty input under TKF92.
We therefore include the singlet term explicitly.

\subsection{Per-branch path log-likelihood}
\label{sec:varanc-presence-path-LL}

Consider a single branch $v \to w$ of length $\branchlen$.
Reading the parent and child rows of the MSA column by column gives,
for each $n \in \{1, \ldots, L\}$, an ordered pair
$(X^v_n, X^w_n) \in \{0,1\}^2$.
Map each pair to a WFST column-state
$S_n \in \{\mat, \ins, \del, \mathsf{Ig}\}$ via
\begin{equation}
\label{eq:state-map}
S_n =
\begin{cases}
  \mat & \text{if } (X^v_n, X^w_n) = (1,1) \\
  \ins & \text{if } (X^v_n, X^w_n) = (0,1) \\
  \del & \text{if } (X^v_n, X^w_n) = (1,0) \\
  \mathsf{Ig} & \text{if } (X^v_n, X^w_n) = (0,0),
\end{cases}
\end{equation}
where the symbol $\mathsf{Ig}$ (``Ignore'') marks columns absent in
both parent and child and so contributing no transition to the WFST
path.
Add boundary sentinels $S_0 = \sta$ and $S_{L+1} = \fin$.
Stripping out columns with $S_n = \mathsf{Ig}$ collapses the sequence
$(S_0, S_1, \ldots, S_L, S_{L+1})$
to a unique path through the WFST whose nontrivial transitions are
between consecutive non-Ignore columns.
The branch log-likelihood is the sum of the corresponding entries of
$\log \tkfwfst'(\insrate, \delrate, \branchlen, \ext)$ from
equation~\eqref{eq:tkf92-wfst}.
Equivalently, summing over each column's incoming transition,
\begin{equation}
\label{eq:branch-LL}
\log P(X^w \mid X^v, \branchlen, \theta)
= \sum_{N=1}^{L+1}
   \delta(S_N \neq \mathsf{Ig})
   \sum_{M=0}^{N-1}
     \delta(S_M \neq \mathsf{Ig})
     \log \tkfwfst'_{S_M S_N}
     \prod_{K=M+1}^{N-1} \delta(S_K = \mathsf{Ig}),
\end{equation}
where $\theta = (\insrate, \delrate, \ext)$,
$\delta(\cdot)$ is a Kronecker indicator,
and the empty product (when $M+1 > N-1$) equals~$1$.
For each non-Ignore column $N$ the inner sum picks out the unique
$M < N$ that is also non-Ignore with all columns strictly between
$M$ and $N$ in state $\mathsf{Ig}$;
that pair contributes $\log \tkfwfst'_{S_M S_N}$ and all other
$M$ contribute zero.

\subsection{Tree-structured variational family}
\label{sec:varanc-presence-q}

We approximate the joint posterior
$P(\{X^v_n\}_{v \in \mathcal{I}} \mid \{X^v_n\}_{v \in \mathcal{L}})$
by a product over MSA columns
\begin{equation}
\label{eq:q-factor}
q(\{X^v_n\}_{v \in \mathcal{I}, 1 \le n \le L})
= \prod_{n=1}^{L} q_n(\{X^v_n\}_{v \in \mathcal{I}}),
\end{equation}
where each per-column distribution $q_n$ is a directed graphical
model on the tree with node values in
\begin{equation}
\label{eq:Z-states}
\mathcal{Z} = \{\textsc{NotYetInserted},\ \textsc{Present},\ \textsc{Deleted}\}
\equiv \{\state{N}, \state{P}, \state{D}\}.
\end{equation}
The variable at the root is drawn from a free distribution
$q_n^{\text{root}} \in \Delta^{\mathcal{Z}}$;
along each edge $v \to w$ the variable $Z^w_n \in \mathcal{Z}$ is
drawn conditionally on $Z^v_n$ from a free distribution
$q_n^{v \to w}(\cdot \mid \cdot)$.
The 9-entry matrix $q_n^{v \to w}(z' \mid z)$ has 4 entries pinned to
zero to enforce irreversibility:
\begin{equation}
\label{eq:irreversibility}
q_n^{v \to w}(\state{D} \mid \state{N}) =
q_n^{v \to w}(\state{N} \mid \state{P}) =
q_n^{v \to w}(\state{N} \mid \state{D}) =
q_n^{v \to w}(\state{P} \mid \state{D}) = 0.
\end{equation}
The two unconstrained rows ($\state{N}$ and $\state{P}$ parents) each
have two surviving outcomes, contributing one free parameter per row;
the $\state{D}$ row is deterministic
($q_n^{v \to w}(\state{D} \mid \state{D}) = 1$).
This leaves two free parameters per (edge, column).
Leaf-node states $Z^v_n$ for $v \in \mathcal{L}$ are clamped: if
$X^v_n = 1$, then $Z^v_n = \state{P}$; if $X^v_n = 0$, then
$Z^v_n \in \{\state{N}, \state{D}\}$ with the variational mass split
between the two by $q_n$.

The presence indicator at internal node $v$ is recovered as
\begin{equation}
\label{eq:Z-to-X}
X^v_n = \delta(Z^v_n = \state{P}),
\end{equation}
so that any realisation drawn from $q_n$ has the property that the
nodes with $X^v_n = 1$ form a connected subtree of $\parsetree$
(or the empty set, which by hypothesis cannot arise once the leaf
clamps are respected, since every column of the supplied MSA contains
at least one Present leaf).

\paragraph{Why this support.}
Under the strict TKF92 process, each MSA column corresponds to a
single insertion event in the history of the tree; the residue
descends, and the set of Present nodes is necessarily a connected
subtree containing the insertion point.
A factorised mean-field $q$ with independent Bernoulli marginals at
each $(v, n)$ would assign positive mass to disconnected
Present-sets, which (i) are unreachable under the model, and
(ii) would force the substitution likelihood to factor as a product
of independent Felsenstein recursions on the two pieces — corresponding
to two separate insertions co-located in one MSA column,
a different generative process.
The 3-state irreversible $q$ rules out both pathologies by
construction.

\subsection{\texorpdfstring{Expected log-likelihood under $q$}{Expected log-likelihood under q}}
\label{sec:varanc-presence-elbo}

For a parent-child pair $(v, w)$ on the tree, define the per-column
state probability under $q_n$:
\begin{align}
\label{eq:column-state-probs}
P_q(S_n = \mat)
  &= q_n(Z^v_n = \state{P},\, Z^w_n = \state{P}), \\
P_q(S_n = \del)
  &= q_n(Z^v_n = \state{P},\, Z^w_n = \state{D}), \\
P_q(S_n = \ins)
  &= q_n(Z^v_n = \state{N},\, Z^w_n = \state{P}), \\
P_q(S_n = \mathsf{Ig})
  &= q_n(Z^v_n = \state{N},\, Z^w_n = \state{N})
   + q_n(Z^v_n = \state{D},\, Z^w_n = \state{D}).
\end{align}
The four other $(z, z')$ joint marginals are zero by
equation~\eqref{eq:irreversibility}.
Each pairwise marginal
$q_n(Z^v_n, Z^w_n)$ is computed by belief propagation
(Section~\ref{sec:varanc-presence-bp}).

The product factorisation~\eqref{eq:q-factor} makes the per-branch
expected log-likelihood factor across columns:
\begin{equation}
\label{eq:E-branch-LL}
\expect_q\!\big[\log P(X^w \mid X^v, \branchlen, \theta)\big]
= \sum_{s, s'}
   \log \tkfwfst'_{s s'}\,
   W_{s s'}^{(v \to w)},
\end{equation}
where $s$ ranges over $\smide \setminus \{\fin\}$ and $s'$ ranges
over $\smide \setminus \{\sta\}$ (so $\sta$ enters only as a source
and $\fin$ only as a destination, both via the deterministic boundary
sentinels), and
\begin{equation}
\label{eq:W-pair}
W_{s s'}^{(v \to w)}
\;=\;
\sum_{N=1}^{L+1}
\sum_{M=0}^{N-1}
P_q(S_M = s)\,P_q(S_N = s')
\prod_{K=M+1}^{N-1} P_q(S_K = \mathsf{Ig})
\end{equation}
is the expected number of times the path uses transition
$s \to s'$ on this branch.
Boundary terms have $P_q(S_0 = \sta) = 1$ and
$P_q(S_{L+1} = \fin) = 1$ deterministically.

\paragraph{Root-prior contribution.}
Under the per-column factorisation~\eqref{eq:q-factor}, the
expected singlet log-likelihood at the root
(equation~\eqref{eq:singlet-tkf92}) reduces to
\begin{equation}
\label{eq:E-singlet-root}
\expect_q\!\big[\log p_{\text{singlet}}(X^{\text{root}})\big]
= \log p \cdot \sum_{n=1}^{L} q_n(X^{\text{root}}_n = 1)
   + \big(1 - P_q(L_{\text{root}}{=}0)\big)\,c_1
   + P_q(L_{\text{root}}{=}0)\,\log(1 - \kappa),
\end{equation}
with $c_1 = \log\!\big[\kappa(1-\ext)(1-\kappa)/p\big]$ and
$P_q(L_{\text{root}}{=}0) = \prod_{n=1}^{L} q_n(X^{\text{root}}_n = 0)$
the variational mass on an empty root sequence.
For TKF91 ($\ext = 0$) the second and third terms collapse to a
single $q$-independent constant $\log(1-\kappa)$ and~\eqref{eq:E-singlet-root}
reduces to the textbook form
$\log(1-\kappa) + \log\kappa \cdot \sum_n q_n(X^{\text{root}}_n = 1)$.
For TKF92 the $L_{\text{root}}{=}0$ special case must be retained;
note that $P_q(L_{\text{root}}{=}0)$ is small but generically nonzero,
since some columns may have positive variational mass on root being
$\state{N}$ (the column was inserted on a strict subtree below the
root).

\subsection{ELBO}
\label{sec:varanc-presence-elbo-derivation}

For any product-of-trees $q$ that gives joint marginals
$q_n(Z^v_n, Z^w_n)$ at every (edge, column),
\begin{equation}
\label{eq:elbo}
\boxed{\;\;
\log p(\text{MSA} \mid \parsetree, \theta)
\;\geq\;
\log \tilde{p}(\text{MSA} \mid \parsetree, \theta)
\;\geq\;
\mathcal{L}[q]
\;=\;
\expect_q\!\big[\log p_{\text{singlet}}(X^{\text{root}})\big]
+ \sum_{(v \to w) \in \parsetree}
   \expect_q\!\big[\log P(X^w \mid X^v, \branchlen_{vw}, \theta)\big]
+ H\!\big[q(\{X^v\}_{v \in \mathcal{I}} \mid \text{MSA})\big].
\;\;}
\end{equation}
The first two terms are given by~\eqref{eq:E-singlet-root}
and~\eqref{eq:E-branch-LL} respectively;
the third is the entropy of the leaf-conditioned variational
distribution.
The leftmost inequality is~\eqref{eq:tildep-le-p};
the right inequality is the standard Jensen bound on the restricted
joint~\eqref{eq:joint}.

\paragraph{Entropy decomposition.}
The variational $q$ is parameterised through a directed graphical
model on the full tree (root + internal + leaf nodes) with edge
conditionals $q^{v \to w}_n(\cdot \mid \cdot)$.
Writing $q_{\text{joint}}$ for this joint and conditioning on the
observed leaf states gives
$q(\{X^v\}_{v \in \mathcal{I}} \mid \text{MSA})
 = q_{\text{joint}}(\text{internals},\,\text{MSA})\,/\,Z_q$,
where $Z_q = \sum_{\text{internals}} q_{\text{joint}}$ is the
per-column belief-propagation partition function summed over columns.
The leaf-conditioned entropy then decomposes as
\begin{equation}
\label{eq:entropy}
H\!\big[q(\text{internals} \mid \text{MSA})\big]
= -\,\expect_q\!\big[\log q_{\text{joint}}\big] + \log Z_q,
\end{equation}
where the first term factors over the directed graphical model:
\begin{equation}
\label{eq:cross-entropy-decomp}
-\expect_q\!\big[\log q_{\text{joint}}\big]
= -\sum_{n=1}^{L}\!\!\sum_{z \in \mathcal{Z}}\!
       q_n^{\text{root}}(z)\,\log q_n^{\text{root}}(z)
\;-\;
\sum_{(v \to w) \in \parsetree}
\sum_{n=1}^{L}\!\!\sum_{z, z' \in \mathcal{Z}}\!
   q_n(Z^v_n{=}z, Z^w_n{=}z')\,\log q_n^{v \to w}(z' \mid z).
\end{equation}
The first term is the per-column root entropy.
Each per-edge term is a \emph{cross-entropy} between the BP-derived
joint marginal $q_n(Z^v_n, Z^w_n)$ and the prior edge conditional
$q_n^{v \to w}$;
it is \emph{not} the parent-marginal-weighted entropy
$\expect_{q_n(Z^v_n)}\!\big[H(q_n^{v \to w}(\cdot \mid Z^v_n))\big]$,
which only equals~\eqref{eq:cross-entropy-decomp} when the leaves
impose no evidence on the edge — generically a different quantity
once the BP-conditioning shifts pair marginals away from the prior
factorisation.
The $\log Z_q$ term in~\eqref{eq:entropy} comes out of the up-pass
partition function and ensures the entropy is on the conditioned
distribution rather than the joint.

\paragraph{Bound interpretation.}
Equation~\eqref{eq:elbo} is a proper lower bound on the full TKF92
marginal log-likelihood $\log p(\text{MSA})$, with the gap split into
two non-negative pieces:
\begin{equation}
\label{eq:gap-decomp}
\log p(\text{MSA}) - \mathcal{L}[q]
\;=\;
\underbrace{\big(\log p(\text{MSA}) - \log \tilde{p}(\text{MSA})\big)}_{\text{restriction gap}}
\;+\;
\underbrace{\text{KL}\!\big[q(\text{internals} \mid \text{MSA})\,\|\,\tilde{p}(\text{internals} \mid \text{MSA})\big]}_{\text{variational gap}}.
\end{equation}
The variational gap is the standard Jensen slack on the restricted
joint~\eqref{eq:joint} and depends on $q$;
the restriction gap is the contribution of ghost-column histories
that our latent space cannot represent
(Section~\ref{sec:varanc-presence-restricted}) and depends on the
model parameters $\theta$ and the tree but not on $q$.
For ancestral-presence reconstruction (fix $\theta$, optimise $q$)
the restriction gap is a constant additive offset, so
maximising~\eqref{eq:elbo} is equivalent to minimising the KL
divergence from $q$ to the true model posterior on the support of
observed-column internal-node configurations.
For parameter learning over $\theta$ this equivalence breaks: the
restriction gap shifts with $\theta$, so the bound is not tight up to
a parameter-independent constant.
A tighter bound would require augmenting $q$ to model ghost columns
explicitly — for instance via per-(branch, gap-region) latent counts
of inserted-then-fully-deleted residues — at the cost of additional
machinery beyond the per-MSA-column factorisation used here.
Equivalently and intuitively, one recovers the full marginal in the
limit of inserting infinitely many additional ``empty'' columns into
the alignment between (and around) every observed column, with all
leaves clamped to absent at those columns and the variational $q$
free over their internal states;
the present formulation simply truncates this padding to zero.

\subsection{Stable computation: cumulant trick}
\label{sec:varanc-presence-cumulant}

The expected transition count $W_{s s'}^{(v \to w)}$
in~\eqref{eq:W-pair} appears at first sight to require
$O(L^2)$ work per (branch, state-pair).
A standard prefix-sum reduces this to $O(L)$.

For brevity drop the branch superscript and write
$P_n^{\mathsf{Ig}} = P_q(S_n = \mathsf{Ig})$,
$P_n^{s} = P_q(S_n = s)$ for $s \in \{\sta, \mat, \del, \ins, \fin\}$
(with $P_0^{\sta} = 1$, $P_{L+1}^{\fin} = 1$, all other boundary
values zero).
Define the running log-Ignore mass
\begin{equation}
\label{eq:cumulant}
C_N \;=\; \sum_{k=1}^{N} \log P_k^{\mathsf{Ig}},
\qquad C_0 = 0.
\end{equation}
Then $\prod_{K=M+1}^{N-1} P_K^{\mathsf{Ig}} = \exp(C_{N-1} - C_M)$,
and~\eqref{eq:W-pair} factorises as
\begin{equation}
\label{eq:W-prefix}
W_{s s'}
\;=\;
\sum_{N=1}^{L+1}
P_N^{s'}\,e^{C_{N-1}}
\;\cdot\;
\underbrace{\sum_{M=0}^{N-1} P_M^{s}\,e^{-C_M}}_{\text{prefix sum in }M}.
\end{equation}
Both inner factors are non-negative; the difference $C_{N-1} - C_M$
is non-positive (it is a sum of log-probabilities), so the geometric
factor lies in $(0, 1]$ and~\eqref{eq:W-prefix} is numerically benign
in float64 for any practical $L$.
For very long alignments where individual exponents could underflow,
one $\epsilon$-floors $P_n^{\mathsf{Ig}}$ before computing $C_n$;
the floor is harmless because the product weights paths that traverse
many Ignore columns, which contribute negligibly to the expectation.

The total work for the per-branch expected log-likelihood
is $O(L \cdot |S|^2)$ with $|S| \le 5$ states,
yielding $O(B L)$ overall for $B$ branches.
This is fully vectorisable across columns and branches.

\paragraph{Why the inner sum is in probability space.}
Equation~\eqref{eq:branch-LL} for a fixed deterministic realisation
of the indicators is purely a sum of $\log \tkfwfst'$ entries:
no probabilities, no products, no exponentials.
The probability-space arithmetic in~\eqref{eq:W-prefix} arises only
when we marginalise over the variational distribution: each summand
is a probability mass (product of independent column factors) times a
real-valued $\log \tkfwfst'$ entry, and the sum collapses these
weighted log-transitions into the expected transition count.

\subsection{Belief propagation for pairwise marginals}
\label{sec:varanc-presence-bp}

The pairwise marginals $q_n(Z^v_n, Z^w_n)$ feeding
equations~\eqref{eq:column-state-probs} are computed by the standard
Felsenstein up-down algorithm on $\parsetree$.
For each column $n$:
\begin{enumerate}[nosep]
\item \emph{Up pass:} for every node $v$ and every state
  $z \in \mathcal{Z}$, accumulate
  $\beta_v^n(z) = \prod_{w: \text{child of } v} \sum_{z'}
                  q_n^{v \to w}(z' \mid z)\,\beta_w^n(z')$,
  with leaf base case
  $\beta_v^n(z) = \delta(\text{leaf clamp at } v, n \text{ is consistent with } z)$.
\item \emph{Down pass:} for every internal $v$ and state $z$,
  accumulate the root-conditioned posterior
  $\alpha_v^n(z)$ in standard fashion.
\item \emph{Marginals:} for every edge $v \to w$ and every state pair,
  $q_n(Z^v_n = z, Z^w_n = z')
    = \alpha_v^n(z)\,q_n^{v \to w}(z' \mid z)\,
      \beta_w^n(z') / \mathcal{Z}_n$,
  where $\mathcal{Z}_n$ is the per-column partition function from the
  up pass at the root.
\end{enumerate}
The up and down passes are $O(|\mathcal{I}| \cdot |\mathcal{Z}|^2)$
per column and trivially vectorise across columns
(same tree topology, different leaf clamps and different per-column
edge conditionals).
For typical protein alignments
($L \approx 360$ amino acids on average in UniProtKB/Swiss-Prot
\cite{UniProtStatistics}, with the long tail extending to a few thousand)
this is unproblematic; for DNA MSAs (gene-scale regions in the few-kbp
range \cite{XuEtAl2006}, longer for syntenic blocks)
the per-column independence keeps cost linear in $L$.

\subsection{Special cases and scalability}
\label{sec:varanc-presence-special-cases}

\paragraph{Irreversible Felsenstein-3.}
Tie the per-(edge, column) variational conditionals
$q_n^{v \to w}(\cdot \mid \cdot)$ across columns to a single
per-edge transition matrix $\bar{q}^{v \to w}(\cdot \mid \cdot)$ —
equivalently, parameterise each $\bar{q}^{v \to w}$ as the
transition matrix of an irreversible 3-state CTMC over
$(\state{N}, \state{P}, \state{D})$ with rates fitted by maximum
likelihood on the leaf data.
The variational marginals $q_n$ are then exactly the Felsenstein
posteriors of this CTMC at column $n$, and the ELBO reduces to the
log-likelihood of the leaf clamps under that CTMC plus the (now
column-independent) expected log of $\tkfwfst'$.

\paragraph{Fitch parsimony.}
The zero-temperature limit of the irreversible Felsenstein-3
posterior — assigning each internal node deterministically to its
most-likely state under the CTMC — recovers the Fitch parsimony
labelling of presence/absence under uniform priors and unit
transition costs.

\paragraph{Scalability of stochastic ELBO optimisation.}
The free-parameter count scales as $2 \cdot |\mathcal{E}| \cdot L$
(plus a per-column root parameter)
where $|\mathcal{E}|$ is the number of tree edges and $L$ is the
number of MSA columns.
The variational $q$ factorises over columns, but the ELBO is
non-trivially coupled across columns through the prefix-of-Ignore
factor in equation~\eqref{eq:W-prefix}, which weights each
$(M, N)$ transition contribution by the chain of intervening Ignore
probabilities.
Match-rich columns (where most lineages are clamped to $\state{P}$)
have $P_n^{\mathsf{Ig}} \approx 0$ and so act as anchors that break
the chain: contributions from $(M, N)$ pairs that straddle such an
anchor column are damped to zero, decoupling the gradient on either
side.
For both protein-scale and DNA-scale alignments with non-trivial
column coverage, ELBO optimisation should therefore scale linearly
with total residue count without running into long-range coupling
pathologies.
